% generated by roam.py handbook export -- не редактировать, править заметки
% 20260922124251-метопт_лекция_1_выпуклые_множества_и_функции.org

\chapter{Выпуклые множества и функции}\label{lecture:1}
\InputIfFileExists{text/01-intro}{}{}

\section{Выпуклые множества}\label{lecture:1:1}  % конспект 1.1
\InputIfFileExists{text/01-1-intro}{}{}

\InputIfFileExists{text/316B5BFA-8B9F-4267-9086-D6593AAB5F9B-before}{}{}
% generated by roam.py handbook export -- не редактировать, править заметки
% 20260916142253-выпуклое_множество.org (body 650a7e662171)

\begin{definition}[{Выпуклое множество}]\label[definition]{note:316B5BFA-8B9F-4267-9086-D6593AAB5F9B}
  Множество выпуклое, если отрезок, соединяющий любые две точки этого множества, лежит полностью в этом множестве.

  Пусть $V$~--- векторное пространство с вещественным полем, $C \subset V$,
  тогда $C$~--- выпуклое, если $\forall x, y \in C\ \forall \lambda \in [0;1]\ \lambda x + (1- \lambda)y \in C$

  Свойства:

  \begin{itemize}
    \item \hyperref[note:08CF7113-4709-4C4F-A0E7-4A2ED595BD7F]{Выпуклость пересечения выпуклых множеств}
    \item \hyperref[note:80AFB48E-12EF-4D6C-9303-E2E52FC83D26]{Выпуклость проекции}
    \item \hyperref[note:88ABE812-A1C1-44A1-BF92-D714C9AF6315]{Выпуклость суммы Минковского выпуклых множеств}
    \item \hyperref[note:42DFF6DC-CC9D-40AA-9D6F-218DCF613833]{Характеристика выпуклого множества через выпуклые комбинации}
  \end{itemize}
\end{definition}
  % Выпуклое множество -- опр. 1.1

\InputIfFileExists{text/08CF7113-4709-4C4F-A0E7-4A2ED595BD7F-before}{}{}
% generated by roam.py handbook export -- не редактировать, править заметки
% 20260916142707-выпуклость_пересечения_выпуклых_множеств.org (body f2f19e21819f)

\begin{theorem}[{Выпуклость пересечения выпуклых множеств}]\label[theorem]{note:08CF7113-4709-4C4F-A0E7-4A2ED595BD7F}
  Пусть $C_{\alpha}$~--- семейство \hyperref[note:316B5BFA-8B9F-4267-9086-D6593AAB5F9B]{выпуклых множеств}, тогда $\bigcap_{\alpha} C_{\alpha}$~--- выпуклое.
\end{theorem}

\begin{proof}
  \[
  x, y \in \bigcap_{\alpha} C_{\alpha} \Rightarrow \forall \alpha\ x,y \in C_{\alpha} \Rightarrow
   \forall \alpha\ \forall 0 \le \lambda \le 1\  \lambda x + (1-\lambda) y \in C_{\alpha} \Rightarrow
  \]
  \[
  \Rightarrow \forall 0 \le \lambda \le 1\ \lambda x + (1-\lambda) y \in \bigcap_{\alpha} C_{\alpha}
  \]
\end{proof}
  % Выпуклость пересечения выпуклых множеств -- свойство 1.1.1

\InputIfFileExists{text/80AFB48E-12EF-4D6C-9303-E2E52FC83D26-before}{}{}
% generated by roam.py handbook export -- не редактировать, править заметки
% 20260916143716-выпуклость_проекции.org (body ba73a7f085ce)

\begin{theorem}[{Выпуклость проекции}]\label[theorem]{note:80AFB48E-12EF-4D6C-9303-E2E52FC83D26}
  Пусть:

  \begin{itemize}
    \item $V$~--- \hyperref[note:5D688904-8445-47B5-9D39-8BCC39F9C354]{гильбертово пространство},
    \item $U \subset V$~--- замкнутое подпространство,
    \item $C \subset V$~--- \hyperref[note:316B5BFA-8B9F-4267-9086-D6593AAB5F9B]{выпуклое множество}
    \item $P_U: V \rightarrow U$~--- проекция (по свойствам гильбертовости это линейный оператор)
  \end{itemize}

  Тогда: $P_U C$~--- тоже выпуклое
\end{theorem}

\begin{proof}
  $x, y \in P_UC$, следовательно у них есть прообразы: $x = P_Ua, y = P_Ub$, где $a,b \in C$.
  Мы хотим доказать, что $\forall \lambda \in [0;1]\ \forall x, y \in P_UC\ \lambda x + (1-\lambda)y \in P_UC$.
  Заменим $x,y$ на проекции прообразов: $\lambda x + (1-\lambda) y = \lambda P_U a + (1- \lambda) P_U b$.
  По линейности: $\lambda P_Ua + (1-\lambda)P_U b = P_U(\lambda a + (1-\lambda) b)$.
  Но мы знаем, что $\lambda a + (1-\lambda)b \in C$ (выпуклость), тогда
  $P_U(\lambda a + (1-\lambda) b) = \lambda P_U a + (1-\lambda) P_U b \in P_U C$
\end{proof}
  % Выпуклость проекции -- свойство 1.1.2

\InputIfFileExists{text/88ABE812-A1C1-44A1-BF92-D714C9AF6315-before}{}{}
% generated by roam.py handbook export -- не редактировать, править заметки
% 20260916213842-выпуклость_суммы_минковского_выпуклых_множеств.org (body d23486e705a1)

\begin{theorem}[{Выпуклость суммы Минковского выпуклых множеств}]\label[theorem]{note:88ABE812-A1C1-44A1-BF92-D714C9AF6315}
  Пусть $C_k,\ k = 1,2,\dots,n$~--- \hyperref[note:316B5BFA-8B9F-4267-9086-D6593AAB5F9B]{выпуклые множества}, $w_k \in \R$~--- веса, причем они все подмножества
  $V$~--- \hyperref[note:1B626AA4-35C8-4C7C-9D3D-D8C26D98F40C]{векторное пространство}.
  Тогда $C = \sum\limits_{k=1}^n w_k C_k = \left\{ \sum\limits_{k=1}^n w_k x_k,\ x_k \in C_k \right\}$~--- \hyperref[note:316B5BFA-8B9F-4267-9086-D6593AAB5F9B]{выпуклое множество}.
  По сути $C$~--- это \hyperref[note:14FE7B45-F304-4C15-8086-3535589FA04A]{Сумма Минковского} для $C_k$, только взвешенная.
\end{theorem}

\begin{proof}
  Нужно доказать:
  \[
  \forall \lambda \in [0;1]\ \forall x, y \in C\ \lambda x +  (1- \lambda )y \in C
  \]

  Знаем, что $x,y$ представимы в виде суммы элементов из множеств $C_k$.
  Тогда:
  \[
  \lambda x + (1-\lambda)y = \lambda \sum\limits_{k=1}^n w_k x_k + (1-\lambda) \sum\limits_{i=1}^n w_i y_i,\quad x_k,y_k \in C_k
  \]
  Тогда объединим слагаемые по индексу в одну сумму
  \[
  \sum\limits_{k=1}^n \left[ \lambda w_kx_k + (1-\lambda) w_ky_k \right]
  \]
  Пусть $z_k:= \lambda x_k + (1-\lambda)y_k$, тогда слагаемые этой суммы имеют вид $w_kz_k$.
  При этом $z_k \in C_k$, так как $C_k$~--- выпуклое.
  Тогда мы представили $\lambda x + (1-\lambda)y$, как $\sum\limits_{k=1}^n w_kz_k, z_{k} \in C_k$.
  А из этого следует, что $\lambda x + (1-\lambda)y \in C$, как элемент \hyperref[note:14FE7B45-F304-4C15-8086-3535589FA04A]{суммы Минковского}
\end{proof}

\begin{example}
  Пусть $C_1$~--- квадрат со стороной 2, а $C_2$~--- круг радиуса 1. Тогда $C_1 + \frac{1}{2} C_2$~--- квадрат со скругленными углами.
\end{example}
  % Выпуклость суммы Минковского выпуклых множеств -- свойство 1.1.3 и пример 1.1 (квадрат + круг)

\InputIfFileExists{text/1AC657F9-2A3D-47D1-A71D-5B813D961260-before}{}{}
% generated by roam.py handbook export -- не редактировать, править заметки
% 20260917205132-выпуклая_оболочка.org (body 23d1471820a4)

\begin{definition}[{Выпуклая оболочка}]\label[definition]{note:1AC657F9-2A3D-47D1-A71D-5B813D961260}
  Пусть $V$~--- векторное пространство, $M \subset V$,
  \[
  \mathrm{conv}(M) = \bigcap_Q \left\{ Q \subset V, Q \text{ -- выпуклое и } M \subset Q \right\}
  \] --- выпуклая оболочка
  Тут логика почти такая же, как с борелевской оболочкой:
  Выпуклая оболочка $M$~--- пересечение всех \hyperref[note:316B5BFA-8B9F-4267-9086-D6593AAB5F9B]{выпуклых множеств}, содержащих $M$.
\end{definition}
  % Выпуклая оболочка -- опр. 1.2

\InputIfFileExists{text/8583F986-A6E8-4BBB-B2FB-BA0E9BBE58BF-before}{}{}
% generated by roam.py handbook export -- не редактировать, править заметки
% 20260917205703-выпуклая_комбинация.org (body 7e161c152eee)

\begin{definition}[{Выпуклая комбинация}]\label[definition]{note:8583F986-A6E8-4BBB-B2FB-BA0E9BBE58BF}
  Пусть $V$~--- \hyperref[note:1B626AA4-35C8-4C7C-9D3D-D8C26D98F40C]{векторное пространство}. $x_k \in V, 1 \le k \le n, \sum\limits_{k=1}^n \lambda_k = 1, \lambda_k \ge 0$ \\
  Тогда $\sum\limits_{k=1}^n \lambda_k x_k$~--- выпуклая комбинация
\end{definition}
  % Выпуклая комбинация -- опр. 1.3

\InputIfFileExists{text/42DFF6DC-CC9D-40AA-9D6F-218DCF613833-before}{}{}
% generated by roam.py handbook export -- не редактировать, править заметки
% 20260917210321-характеристика_выпуклого_множества_через_выпуклые_комбинации.org (body aa80ca90b346)

\begin{theorem}[{Характеристика выпуклого множества через выпуклые комбинации}]\label[theorem]{note:42DFF6DC-CC9D-40AA-9D6F-218DCF613833}
  Пусть $V$~--- векторное пространство. $C \subset V$.
  $C$~--- \hyperref[note:316B5BFA-8B9F-4267-9086-D6593AAB5F9B]{выпуклое} тогда и только тогда, когда оно содержит все свои \hyperref[note:8583F986-A6E8-4BBB-B2FB-BA0E9BBE58BF]{выпуклые комбинации}.
\end{theorem}

\begin{proof}
  \proofpart{$\Leftarrow$ (Содержит комбинации, следовательно выпуклое)}
  Если $\lambda \in [0;1]$, то частным случаем ($n=2, \lambda_1=\lambda, \lambda_2 = 1-\lambda, x_1=x, x_2 = y$) \hyperref[note:8583F986-A6E8-4BBB-B2FB-BA0E9BBE58BF]{выпуклой комбинации} является $\lambda x + (1-\lambda)y\ \forall x, y \in C$.
  А это в точности определение \hyperref[note:316B5BFA-8B9F-4267-9086-D6593AAB5F9B]{выпуклого множества}.

  \proofpart{$\Rightarrow$ (Выпуклое, следовательно содержит комбинации)}
  Строим индукцию.
  База это определение \hyperref[note:316B5BFA-8B9F-4267-9086-D6593AAB5F9B]{выпуклого множества}.

  Положим, что для $\forall n \le m$ утверждение  из заголовка верно.
  Увеличим размерность суммы.

  Давайте запишем произвольную \hyperref[note:8583F986-A6E8-4BBB-B2FB-BA0E9BBE58BF]{выпуклую комбинацию}
  $\sum\limits_{k=1}^{m+1} \lambda_k x_k,\ x_k \in C, \sum\limits_{k=1}^{m+1} \lambda_k = 1$. Если $\lambda_{m+1} = 1$, то остальные $\lambda_k = 0$ и комбинация равна $x_{m+1} \in C$. Далее считаем $\lambda_{m+1} \ne 1$.

  Тогда:
  $\sum\limits_{k=1}^{m+1} \lambda_k x_k = \sum\limits_{k=1}^{m} \lambda_k x_k + \lambda_{m+1}x_{m+1}$

  Пусть $y = \sum\limits_{k=1}^m \frac{\lambda_k x_k}{1 - \lambda_{m+1}}$. Но $\sum\limits_{k=1}^m \frac{\lambda_k}{1-\lambda_{m+1}} = \frac{1-\lambda_{m+1}}{1-\lambda_{m+1}} = 1$, и коэффициенты неотрицательны ($1 - \lambda_{m+1} > 0$), следовательно $y$~--- \hyperref[note:8583F986-A6E8-4BBB-B2FB-BA0E9BBE58BF]{выпуклая комбинация} и можем применить индукционное предположение для размерности $m$.
  Тогда $y \in C$.

  Получили, что $\sum\limits_{k=1}^{m+1} \lambda_k x_k = \sum\limits_{k=1}^{m} \lambda_k x_k + \lambda_{m+1}x_{m+1} = (1-\lambda_{m+1})y + \lambda_{m+1}x_{m+1}$
  А по \hyperref[note:316B5BFA-8B9F-4267-9086-D6593AAB5F9B]{определению выпуклости} этот объект лежит в множестве.
\end{proof}
  % Характеристика выпуклого множества через выпуклые комбинации -- лемма 1.1

\InputIfFileExists{text/3A5B0E29-AED4-409E-A823-956F3A3957F7-before}{}{}
% generated by roam.py handbook export -- не редактировать, править заметки
% 20260918102313-внутреннее_описание_выпуклои_оболочки.org (body b16acfedea0a)

\begin{theorem}[{Внутреннее описание выпуклой оболочки}]\label[theorem]{note:3A5B0E29-AED4-409E-A823-956F3A3957F7}
  Пусть $V$~--- \hyperref[note:1B626AA4-35C8-4C7C-9D3D-D8C26D98F40C]{векторное пространство}, $M \subset V$.
  Тогда
  \[
  \mathrm{conv}(M) = \left\{ \sum\limits_{k=1}^n \lambda_k x_k\  \left| \  x_k \in M, \sum\limits_{k=1}^n \lambda_k = 1, \lambda_k \ge 0\right. \right\}
  \]

  То есть \hyperref[note:1AC657F9-2A3D-47D1-A71D-5B813D961260]{выпуклая оболочка} множества~--- это всевозможные \hyperref[note:8583F986-A6E8-4BBB-B2FB-BA0E9BBE58BF]{выпуклые комбинации} его элементов.
\end{theorem}

\begin{proof}
  Пусть
  \[
  S = \left\{ \sum\limits_{k=1}^n \lambda_k x_k\  \left| \  x_k \in M, \sum\limits_{k=1}^n \lambda_k = 1, \lambda_k \ge 0\right. \right\}
  \]

  \proofpart{Сперва докажем, что $\mathrm{conv}(M) \supset S$}
  По определению \hyperref[note:1AC657F9-2A3D-47D1-A71D-5B813D961260]{выпуклой оболочки}:
  \[
  \mathrm{conv}(M) = \bigcap_Q \left\{ Q \subset V, Q \text{ -- выпуклое и } M \subset Q \right\}
  \]

  А из \hyperref[note:42DFF6DC-CC9D-40AA-9D6F-218DCF613833]{леммы} мы знаем, что \hyperref[note:316B5BFA-8B9F-4267-9086-D6593AAB5F9B]{выпуклое множество} содержит все свои выпуклые комбинации.

  Получается, что $x_k \in M \Rightarrow x_k \in Q \supset M$, для произвольного выпуклого $Q \supset M$.
  Тогда \hyperref[note:8583F986-A6E8-4BBB-B2FB-BA0E9BBE58BF]{выпуклая комбинация} $x_k \in M$, тоже принадлежит к $Q \supset M$.

  Получили, что $S \subset Q$ для любого выпуклого $Q \supset M$, следовательно $S$ лежит и в их пересечении $\Rightarrow S \subset \mathrm{conv}(M)$

  \proofpart{Теперь докажем, что $S \supset \mathrm{conv}(M)$}
  Так как $\conv(M)$ это пересечение всех \hyperref[note:316B5BFA-8B9F-4267-9086-D6593AAB5F9B]{выпуклых надмножеств} $M$, то для любого
  \hyperref[note:316B5BFA-8B9F-4267-9086-D6593AAB5F9B]{выпуклого} $Q \supset M$ верно, что $Q \supset \conv(M)$.

  Следовательно, нам нужно доказать, что $S$~--- \hyperref[note:316B5BFA-8B9F-4267-9086-D6593AAB5F9B]{выпукло} и $M \subset S$.

  \proofstep{$M \subset S$}
  Так как $S$~--- все \hyperref[note:8583F986-A6E8-4BBB-B2FB-BA0E9BBE58BF]{выпуклые комбинации} $M$, то можем рассмотреть вырожденный случай \hyperref[note:8583F986-A6E8-4BBB-B2FB-BA0E9BBE58BF]{выпуклой комбинации}~--- суммы из одного элемента с $\lambda =1$.

  Тогда получим, что каждый элемент из $M$ лежит и в $S$.

  \proofstep{$S$~--- выпуклое}
  Проверим, по определению \hyperref[note:316B5BFA-8B9F-4267-9086-D6593AAB5F9B]{выпуклого множества}.

  Рассмотрим два элемента $u,v \in S$. Мы знаем, как они устроены (по определению они \hyperref[note:8583F986-A6E8-4BBB-B2FB-BA0E9BBE58BF]{выпуклые комбинации} элементов из $M$):
  \[
  u = \sum\limits_{i=1}^n \lambda_i x_i,\quad v = \sum\limits_{j=1}^m \eta_jy_j,\qquad
  \sum\limits_{i=1}^n \lambda_i = \sum\limits_{j=1}^m \eta_j = 1,\ \eta_j, \lambda_i \ge 0,\ x_i,y_j \in M
  \]

  Возьмем $\phi \in [0;1]$, рассмотрим $\phi u + (1-\phi) v$, подставим суммы
  \[
   \sum\limits_{i=1}^n (\phi \lambda_i) x_i + \sum\limits_{j=1}^m ((1-\phi)\eta_j)y_j
  \]

  Получили комбинацию из $n+m$ точек с коэффициентами $\phi\lambda_i$ и $(1-\phi)\eta_j$ соответственно.

  Коэффициенты неотрицательны; если их сумма 1~--- то это будет \hyperref[note:8583F986-A6E8-4BBB-B2FB-BA0E9BBE58BF]{выпуклой комбинацией}, проверим
  \[
  \sum\limits_{i=1}^n \phi \lambda_i + \sum\limits_{j=1}^m (1-\phi) \eta_j = \phi \cdot \sum\limits_{i=1}^n \lambda_i + (1-\phi) \sum\limits_{j=1}^m \eta_j = \phi \cdot 1 + (1-\phi) \cdot 1 = 1
  \]

  Получили, что $\sum\limits_{i=1}^n (\phi \lambda_i) x_i + \sum\limits_{j=1}^m ((1-\phi)\eta_j)y_j$~--- \hyperref[note:8583F986-A6E8-4BBB-B2FB-BA0E9BBE58BF]{выпуклая комбинация} элементов из $M$, следовательно она лежит в $S$.

  Следовательно, мы доказали, что $\forall u, v \in S\ \forall \phi \in [0; 1]\ \phi u + (1-\phi) v \in S$ \\
  Это в точности определение \hyperref[note:316B5BFA-8B9F-4267-9086-D6593AAB5F9B]{выпуклого множества}
\end{proof}
  % Внутреннее описание выпуклой оболочки -- теорема 1.1

\section{Выпуклые функции}\label{lecture:1:2}  % конспект 1.2
\InputIfFileExists{text/01-2-intro}{}{}

\InputIfFileExists{text/1CA14F2C-637F-41DE-9E97-72075D19E97F-before}{}{}
% generated by roam.py handbook export -- не редактировать, править заметки
% 20260918131811-выпуклая_функция.org (body 061c43aff709)

\begin{definition}[{Выпуклая функция}]\label[definition]{note:1CA14F2C-637F-41DE-9E97-72075D19E97F}
  Пусть $E$~--- \hyperref[note:316B5BFA-8B9F-4267-9086-D6593AAB5F9B]{выпуклое множество}; $f: E \rightarrow \R$.
  $f$~--- выпуклая, если:
  \[\forall x,y \in E\ \forall \lambda \in [0;1]\ f(\lambda x + (1-\lambda)y) \le \lambda f(x) + (1-\lambda)f(y)\]

  Критерии и свойства:

  \begin{itemize}
    \item \hyperref[note:18BD8497-CEBD-4962-8144-474B489606C3]{Дифференциальный критерий выпуклости}
    \item \hyperref[note:86FCE16B-C82E-4416-851C-0E0AD0FDA382]{Критерий выпуклости в терминах вторых производных}
    \item \hyperref[note:1CD538BF-6F40-472A-A670-D38F093E4DF1]{Выпуклость множества подуровня выпуклой функции}
    \item \hyperref[note:035B8935-65A9-4BB4-BDA3-60A60FF278F0]{Наименьшее значение выпуклой функции достигается в точке с нулевым градиентом}
    \item \hyperref[note:6BC7D685-C429-45A3-B563-269840CA0664]{Преобразования, сохраняющие выпуклость}
  \end{itemize}
\end{definition}
  % Выпуклая функция -- опр. 1.4

\InputIfFileExists{text/18BD8497-CEBD-4962-8144-474B489606C3-before}{}{}
% generated by roam.py handbook export -- не редактировать, править заметки
% 20260918132958-дифференциальныи_критерии_выпуклости.org (body c847a4d12c2c)

\begin{theorem}[{Дифференциальный критерий выпуклости}]\label[theorem]{note:18BD8497-CEBD-4962-8144-474B489606C3}
  Пусть $E \subset \R^n$~--- \hyperref[note:316B5BFA-8B9F-4267-9086-D6593AAB5F9B]{выпуклое множество}, $f: E \rightarrow \R$~--- \hyperref[note:957FA09D-8883-40A5-B4C7-237D9F815A42]{дифференцируемая} функция.
  Тогда $f$~--- \hyperref[note:1CA14F2C-637F-41DE-9E97-72075D19E97F]{выпуклая} $\Leftrightarrow$ $\forall x,y \in E\ f(y) \ge f(x) + \left\langle \nabla f(x); y-x \right\rangle$
\end{theorem}

\begin{remark}
  Это аналог \hyperref[note:2F475EE4-2630-45FD-8DD2-8EA9D48059E4]{одномерного случая}, когда \hyperref[note:1CA14F2C-637F-41DE-9E97-72075D19E97F]{выпуклая функция} должна лежать не ниже касательной (в данном случае $f: E \subset \R \rightarrow \R$ и $f$~--- \hyperref[note:957FA09D-8883-40A5-B4C7-237D9F815A42]{дифференцируемая})
  $f(y) \ge f(x) + f^{\prime}(x)(y-x) \Leftrightarrow f$~--- выпуклая функция
\end{remark}

\begin{proof}
  \proofpart{$\Rightarrow$}
  Давайте сведём наш случай к \hyperref[note:2F475EE4-2630-45FD-8DD2-8EA9D48059E4]{одномерному}.
  Так как нас интересует только отрезок $[x;y]$, мы можем сузить функцию $f$ на этот отрезок.
  Заведем $g(t) = f(x + t(y-x)) = f((1-t)x + ty),\ t \in [0;1]$

  \proofstep{$g$~--- выпуклая}
  Возьмем $l,t,\lambda \in [0;1]$, проверим определение выпуклости.
  \[
  g(\lambda l + (1-\lambda) t) = f(x + (\lambda l + (1-\lambda)t)(y-x)) =
  \]
  Представим $x$ как $\lambda x + (1-\lambda) x$:
  \[
  f(\lambda x + (1-\lambda) x + (\lambda l + (1-\lambda)t)(y-x)) =
  \]
  Перегруппируем на слагаемые с $\lambda$ и с $1-\lambda$
  \[
  f (\lambda ( x + l(y-x)) + (1-\lambda)(x + t(y-x)))
  \]
  По выпуклости, получаем неравенство
  \begin{multline*}
  f (\lambda ( x + l(y-x)) + (1-\lambda)(x + t(y-x))) \le \\
  \le \lambda f(x + l(y-x)) + (1-\lambda) f(x + t(y-x)) \qquad\qquad
  \end{multline*}
  Давайте выразим $f$, через $g$.
  \[
  \lambda f(x + l(y-x)) + (1-\lambda) f(x + t(y-x)) = \lambda g(l) + (1-\lambda) g(t)
  \]
  Итого $g$~--- \hyperref[note:1CA14F2C-637F-41DE-9E97-72075D19E97F]{выпуклая}.

  \proofstep{Применим одномерный случай}
  Теперь применим \hyperref[note:2F475EE4-2630-45FD-8DD2-8EA9D48059E4]{одномерный случай} к выпуклой $g$ в точках $0$ и $1$:
  $f(y)= g(1) \ge g(0) + g^{\prime}(0)(1-0) = f(x) + g^{\prime}(0)$ (*)

  Давайте посмотрим на $g^{\prime}(t)$, можно представить $g$, как $f \circ h$, где
  $h(t) = x + t(y-x)$, тогда $g^{\prime}(t) = f^{\prime}(h(t))h^{\prime}(t)$

  Для $h: \R \rightarrow \R^n$ производной является  столбец $n \times 1$. Можно заметить, что она равна множителю при $t$:  $(y-x)$.  $y,x$~--- это векторы, так что получится столбец.
  Для $f: \R^n \rightarrow \R$ производной является строка $1 \times n$

  Тогда, выходит, что $f^{\prime}(h(t)) h^{\prime}(t) = (\nabla f(h(t)))^T h^{\prime}(t)$, где $\nabla f$~--- \hyperref[note:1756907B-1FF2-4F62-8553-759B8509281C]{градиент} \\
  Получилось \hyperref[note:F2942E98-6115-4600-AA3F-E0F687D9C732]{скалярное произведение}; Давайте подставим значение производной и функции
  \[
  (\nabla f(h(t)))^T h^{\prime}(t) =\left\langle \nabla f(x + t(y-x)); y-x \right\rangle
  \]

  В точке $t=0, g^{\prime}(0) = \left\langle \nabla f(x); y-x \right\rangle$

  Подставим это в (*) и получим
  \[
  f(y) \ge f(x) + \left\langle \nabla f(x); y-x \right\rangle
  \]

  \proofpart{$\Leftarrow$}
  Имеем, что $E \subset \R^n$~--- выпуклое, $f : E \rightarrow \R$~--- дифференцируемая и
  \[
  \forall x, y \in E\ f(y) \ge f(x) + \left\langle \nabla f(x); y-x \right\rangle
  \]

  Давайте возьмем произвольные $x,y \in E$, произвольный $\lambda \in [0;1]$  и $z:= \lambda y + (1-\lambda) x$. Мы знаем, что $z \in E$, так как $E$~--- выпуклое.

  Запишем неравенство два раза: для $x, z$ и  для $y,z$

  \[
  f(y) \ge f(z) + \left\langle \nabla f(z); y - z \right\rangle \tag{1}
  \]
  \[
  f(x) \ge f(z) + \left\langle \nabla f(z); x - z \right\rangle \tag{2}
  \]

  Домножим (1) на $\lambda$, а (2) на $1-\lambda$ и сложим их. Получим:
  \[
  \lambda f(y) + (1-\lambda) f(x) \ge \lambda f(z) + (1-\lambda) f(z) + \lambda \left\langle \nabla f(z); y-z \right\rangle + (1-\lambda) \left\langle  \nabla f(z); x-z \right\rangle
  \]
  Давайте занесем множители в скалярное произведение, а также объединим $f(z)$.
  \[
  \lambda f(y) + (1-\lambda) f(x) \ge f(z)+  \left\langle\nabla f(z);  \lambda \cdot ( y-z ) \right\rangle +\left\langle   \nabla f(z); (1-\lambda) \cdot (  x-z ) \right\rangle
  \]
  Давайте теперь подставим значение $z$
  \[
  \lambda f(y) + (1-\lambda) f(x) \ge  f(z)+  \left\langle \nabla f(z); \lambda (  y - \lambda y - (1-\lambda) x )\right\rangle +\left\langle   \nabla f(z); (1-\lambda) ( x - \lambda y - (1-\lambda) x ) \right\rangle
  \]
  Сократим слагаемые с $x, y$
  \[
  \lambda f(y) + (1-\lambda) f(x) \ge  f(z)+  \left\langle \nabla f(z); \lambda ( (1-\lambda) y - (1-\lambda) x )\right\rangle +\left\langle   \nabla f(z); (1-\lambda) ( \lambda x - \lambda y ) \right\rangle
  \]
  Вынесем множители за скобки
  \[
   \lambda f(y) + (1-\lambda) f(x) \ge  f(z)+  \left\langle \nabla f(z); \lambda (1-\lambda) (  y -  x )\right\rangle +\left\langle   \nabla f(z); (1-\lambda) \lambda ( x -y ) \right\rangle
  \]
  По линейности сложим скалярное
  \[
   \lambda f(y) + (1-\lambda) f(x) \ge  f(z)+  \left\langle \nabla f(z); \lambda (1-\lambda) (  y -  x ) + \lambda (1-\lambda) (x - y)\right\rangle
  \]
  Получили
  \[
   \lambda f(y) + (1-\lambda) f(x) \ge  f(z)+  \left\langle \nabla f(z); 0 \right\rangle
  \]
  Подставим $z$ в $f(z)$ и запишем итоговое неравенство.
  Итого  $\forall x,y \in E\ \forall \lambda \in [0;1]\ \lambda f(y) + (1-\lambda) f(x) \ge f(\lambda y + (1-\lambda) x)$ \\
  Получили определение выпуклой функции
\end{proof}
  % Дифференциальный критерий выпуклости -- теорема 1.2 (замечание 1.1 о выпуклости $g(t) = f(x + t(y-x))$ — внутри доказательства)

\InputIfFileExists{text/035B8935-65A9-4BB4-BDA3-60A60FF278F0-before}{}{}
% generated by roam.py handbook export -- не редактировать, править заметки
% 20260918142926-наименьшее_значение_выпуклои_функции_достигается_в_точке_с_нулевым_градиентом.org (body 0d71ae19db42)

\begin{theorem}[{Наименьшее значение выпуклой функции достигается в точке с нулевым градиентом}]\label[theorem]{note:035B8935-65A9-4BB4-BDA3-60A60FF278F0}
  Пусть $E \subset \R^n$~--- \hyperref[note:316B5BFA-8B9F-4267-9086-D6593AAB5F9B]{выпуклое множество}, $f: E \rightarrow \R$~--- \hyperref[note:1CA14F2C-637F-41DE-9E97-72075D19E97F]{выпуклая} и дифференцируемая;
  возьмем  $x^{*} \in E$
  Тогда $\nabla f(x^{*}) = 0 \Leftrightarrow x^{*}$~--- точка минимума ($\nabla f$~--- \hyperref[note:1756907B-1FF2-4F62-8553-759B8509281C]{градиент})
\end{theorem}

\begin{remark}
  Для $\Leftarrow$ нужно, чтобы $x^{*} \in \operatorname{Int} E$
\end{remark}

\begin{proof}
  \proofpart{$\Rightarrow$}
  Возьмем неравенство из \hyperref[note:18BD8497-CEBD-4962-8144-474B489606C3]{дифференциального критерия выпуклости}
  \[
  f(y) \ge f(x) + \left\langle \nabla f(x); y-x \right\rangle\ \forall x, y \in E
  \]
  Если подставить $x^{*}$, вместо $x$, получим
  \[
  \forall y \in E\ f(y) \ge f(x^{*}) + \left\langle \nabla f(x^{*}); y - x^{*} \right\rangle
  \Rightarrow f(y) \ge f(x^{*}) + \left\langle 0; y-x^{*} \right\rangle
  \]
  \[
  \Rightarrow f(y) \ge f(x^{*})\ \forall y \in E
  \]

  Причем оно достигается, так как $x^{*} \in E$.

  \proofpart{$\Leftarrow$}
  Давайте возьмем произвольное направление $h\in \R^n$,
  Так как $x^{*}$~--- внутренняя точка $E$,  найдется окрестность этой точки~--- $U$, такая, что $U \subset E$.
  А из этого следует, что $\exists \varepsilon >0:\ x^{*} + th \in U,\ |t| < \varepsilon$
  Рассмотрим $g(t) = f(x^{*} + th)$, она определена на $(-\varepsilon; \varepsilon)$.
  Так как $x^{*}$~--- точка минимума $f$, то $0$~--- точка минимума $g$, так как $g(0) = f(x^{*})$
  $g$~--- дифференцируема. Тогда по \hyperref[note:1E03F05F-22CC-4D55-BCF1-EBB5774765F0]{теореме Ферма} $g^{\prime}(0) = 0$.

  Что представляет из себя $g^{\prime}(t)$:
  \[
  g^{\prime}(t) = \frac{\partial }{\partial t} f (x^{*} + th) = f^{\prime}(x^{*} + th) h = (\nabla f( x^{*} + th))^T h = \left\langle
  \nabla f(x^{*} +th); h
  \right\rangle
  \]

  Но мы доказали, что для произвольного  $h$, $g^{\prime}(0) = \left\langle \nabla f (x^{*}); h \right\rangle = 0$
  Получается, если взять $h = \nabla f(x^{*})$ и подставить в производную, то
  \[
  \left\langle \nabla f(x^{*}) ; \nabla f(x^{*}) \right\rangle = \| \nabla f (x^{*}) \| ^2 = 0 \Rightarrow \nabla f(x^{*}) = \overline{0}
  \]
\end{proof}
  % Наименьшее значение выпуклой функции достигается в точке с нулевым градиентом -- следствие 1.1

\InputIfFileExists{text/BF27A4AF-11BF-4A0D-85DF-5F7790BD58A9-before}{}{}
% generated by roam.py handbook export -- не редактировать, править заметки
% 20260918155031-линия_уровня_функции.org (body 6aec83aed721)

\begin{definition}[{Линия уровня функции}]\label[definition]{note:BF27A4AF-11BF-4A0D-85DF-5F7790BD58A9}
  Пусть $E$~--- произвольное множество, $f: E \rightarrow \R$, $c \in \R$ \\
  Тогда множество $f_c = \left\{ x \in E \ | \ f(x) = c  \right\}$~--- \emph{линия уровня}
\end{definition}
  % Линия уровня функции -- опр. 1.5

\InputIfFileExists{text/1CD538BF-6F40-472A-A670-D38F093E4DF1-before}{}{}
% generated by roam.py handbook export -- не редактировать, править заметки
% 20260918155412-выпуклость_множества_подуровня_выпуклои_функции.org (body 08743327dc0c)

\begin{theorem}[{Выпуклость множества подуровня выпуклой функции}]\label[theorem]{note:1CD538BF-6F40-472A-A670-D38F093E4DF1}
  Пусть $f: E \rightarrow \R$~--- \hyperref[note:1CA14F2C-637F-41DE-9E97-72075D19E97F]{выпуклая функция}, $c \in \R$, $E$~--- \hyperref[note:316B5BFA-8B9F-4267-9086-D6593AAB5F9B]{выпуклое}
  Тогда $\mathrm{Lev}_f(c) = \left\{ x \in E \  |\  f(x) \le c  \right\}$~--- \hyperref[note:316B5BFA-8B9F-4267-9086-D6593AAB5F9B]{выпуклое множество}
\end{theorem}

\begin{proof}
  Будем проверять определение выпуклого множества.

  Возьмем произвольные $x, y \in \mathrm{Lev}_f(c), \lambda \in [0;1]$.
  Заметим, что $\mathrm{Lev}_f(c) \subset E \Rightarrow x,y \in E$
  Знаем, что $\lambda x + (1-\lambda) y \in E$ \hyperref[note:316B5BFA-8B9F-4267-9086-D6593AAB5F9B]{(по выпуклости)}

  По выпуклости $f$: $f(\lambda x + (1-\lambda)y) \le \lambda f(x) + (1-\lambda)f(y)$

  \[
  x,y \in \mathrm{Lev}_f(c) \Rightarrow \lambda f(x) + (1-\lambda)f(y) \le \lambda c + (1-\lambda)c = c
  \]

  Но тогда, $f(\lambda x + (1-\lambda)y) \le c \Rightarrow (\lambda x + (1-\lambda)y) \in \mathrm{Lev}_{f}(c)$
  Итого $\mathrm{Lev}_f(c)$~--- выпуклое \hyperref[note:316B5BFA-8B9F-4267-9086-D6593AAB5F9B]{по определению}.
\end{proof}
  % Выпуклость множества подуровня выпуклой функции -- утверждение 1.1 ($\mathrm{Lev}_f(c)$)

\InputIfFileExists{text/D053F5F9-AF9F-4517-A434-44B68A40039F-before}{}{}
% generated by roam.py handbook export -- не редактировать, править заметки
% 20260918160713-выпуклость_линии_уровня_для_инфимума_функции.org (body 3cd405f5e09a)

\begin{theorem}[{Выпуклость линии уровня для инфимума функции}]\label[theorem]{note:D053F5F9-AF9F-4517-A434-44B68A40039F}
  Если $f: E \rightarrow \R$~--- \hyperref[note:1CA14F2C-637F-41DE-9E97-72075D19E97F]{выпуклая}, $m = \inf_{x \in E} f(x)$
  Тогда $f_m$ (\hyperref[note:BF27A4AF-11BF-4A0D-85DF-5F7790BD58A9]{линия уровня})~--- \hyperref[note:316B5BFA-8B9F-4267-9086-D6593AAB5F9B]{выпуклое множество}.
\end{theorem}

\begin{remark}
  Если инфимум не достигается, то $f_m = \varnothing$~--- тоже выпуклое множество.
\end{remark}

\begin{proof}
  Мы знаем, что множество подуровня выпуклой функции~--- выпукло (см. \hyperref[note:1CD538BF-6F40-472A-A670-D38F093E4DF1]{Выпуклость множества подуровня выпуклой функции})

  А так как $\mathrm{Lev}_f(c) = \left\{ x \in E\ |\ f(x) \le c  \right\}$, и $f(x) \ge m\ \forall x$ ($m$~--- инфимум)
  Получим
  \[
  \mathrm{Lev}_f(m) = \left\{ x  \in E\ |\ f(x) \le m \right\} = \left\{ x \in E\ |\ f(x) = m \right\} = f_m
  \]
  Получили в точности \hyperref[note:BF27A4AF-11BF-4A0D-85DF-5F7790BD58A9]{линию уровня функции}
\end{proof}
  % Выпуклость линии уровня для инфимума функции -- следствие 1.2

\InputIfFileExists{text/86FCE16B-C82E-4416-851C-0E0AD0FDA382-before}{}{}
% generated by roam.py handbook export -- не редактировать, править заметки
% 20260918164233-критерии_выпуклости_в_терминах_вторых_производных.org (body d21a2f5fd79d)

\begin{theorem}[{Критерий выпуклости в терминах вторых производных}]\label[theorem]{note:86FCE16B-C82E-4416-851C-0E0AD0FDA382}
  Пусть $E \subset \R^n$~--- открытое выпуклое множество, $f: E \rightarrow \R$~--- дважды дифференцируемая;

  Тогда $f$~--- выпуклая $\Leftrightarrow$ $\nabla^2 f(x) \ge 0\ \forall x \in E$
\end{theorem}

\begin{remark}
  Правое условие аналогично $\lambda_{\min}(\nabla^2 f(x)) \ge 0$~--- \hyperref[note:025A8419-962C-4892-A4F8-FB496AAA938C]{факт из алгебры}
\end{remark}

\begin{proof}
  Рассмотрим для произвольных $x \in E$, $h \in \R^n$ функцию $\phi(t) = f(x + th)$

  \proofpart{Выпуклость $\phi \Leftrightarrow$ выпуклости $f$}
  \proofstep{$\Leftarrow$}
  Докажем, что $\phi$~--- выпуклая. Зададим область определения $\phi$
  \[
  I = \left\{ t \in \R\ |\ x + th \in E \right\}
  \]

  Заметим, что $I$~--- выпукло, так как для любых $t_1, t_2 \in I$
  точки $x + t_1h \in E$  и $x+ t_2h \in E$ \\
  Но $E$~--- выпукло, следовательно отрезок $[x + t_1 h; x + t_2h] \subset E$,
  а из этого следует, что отрезок $[t_1; t_2] \subset I$, так как $\forall t \in [t_1; t_2]\ x + th \in E$ \\
  А это ровно так, как мы определяем принадлежность к $I$.

  Возьмем произвольные $l, t \in I$ и $\lambda \in [0;1]$
  \[
  \phi(\lambda t + (1-\lambda) l) = f(x + (\lambda t + (1-\lambda)l)h)
  \]
  Представим $x = \lambda x + (1-\lambda) x$ и сгруппируем вокруг $\lambda, 1-\lambda$
  \[
  \phi(\lambda t + (1-\lambda) l) = f(\lambda (x + th) + (1-\lambda)(x + lh))
  \]

  Так как $l, t \in I \Rightarrow (x + th), (x+lh) \in E$, тогда из выпуклости $E$ следует, что
  $\lambda (x + th) + (1-\lambda)(x + lh) \in E$.

  Тогда, можно записать неравенство:
  \[
  f(\lambda (x + th) + (1-\lambda)(x + lh)) \le \lambda f(x + th) + (1-\lambda)f(x + lh)
  \]

  Но $\lambda f(x + th) + (1-\lambda)f(x + lh) = \lambda \phi(t) + (1-\lambda) \phi(l)$ \\
  Тогда мы получим $\phi(\lambda t + (1-\lambda) l) \le \lambda \phi(t) + (1 - \lambda)\phi(l)\ \forall l,t \in I, \lambda \in [0;1]$ \\
  Мы доказали, что $\phi$~--- выпуклая.

  \proofstep{$\Rightarrow$}
  Давайте возьмем произвольные $x, y \in E,\ h := y-x$.
  Тогда точки $x + th$ при $t \in [0;1]$ пробегают отрезок $[x;y]$, а $\phi(t) := f(x + th)$~--- сужение $f$ на эту прямую; по предположению $\phi$~--- выпуклая.
  Давайте будем проверять $f$ на выпуклость.
  \[
  f(\lambda y + (1-\lambda) x) = f (x + \lambda (y - x)) = \phi(\lambda) \le \lambda \phi(1) + (1-\lambda) \phi(0) = \lambda f(y) + (1-\lambda) f(x)
  \]

  \proofpart{Связь выпуклости $\phi$ и $\nabla^2 f(x) \ge 0$}
  \proofstep{$\Rightarrow$}
  Мы все еще пользуемся произвольными $x \in E$, $h \in \R^n$.

  Рассмотрим $\phi^{\prime}(t) = \left( f(x+th) \right)^{\prime}_t = f^{\prime}(x + th) (x + th)^{\prime}_t = \left\langle \nabla f (x + th); h \right\rangle$ (см. \hyperref[note:1756907B-1FF2-4F62-8553-759B8509281C]{градиент}) \\
  Вторая производная:
  \[
  \phi^{\prime \prime}(t) = \left(  \left\langle \nabla f(x +th); h \right\rangle \right)^{\prime}_t =
  \left\langle \left( \nabla f(x +th) \right)^{\prime}_t; h \right\rangle  +
   \left\langle \nabla f(x +th); h^{\prime}_t \right\rangle
  \]
  Получим
  \[
  \phi^{\prime \prime}(t) = \left\langle \nabla^2 f(x +th) h; h \right\rangle + \left\langle \nabla f(x+th); 0 \right\rangle
  = \left\langle \nabla^2 f(x +th) h; h \right\rangle
  \]
  Раскроем скалярное произведение
  \[
  \left\langle \nabla^2 f(x +th) h; h \right\rangle = h^T \nabla^2 f(x +th) h
  \]

  Из \hyperref[note:9CEE854D-ED37-4C08-82EB-E41B36C43EA6]{одномерного случая}, мы знаем, что $\phi$~--- выпуклая $\Leftrightarrow \phi^{\prime\prime}(t) \ge 0\ \forall t \in I$.
  В частности мы имеем, что $\phi^{\prime \prime} (0) = h^T \nabla^2 f(x) h \ge 0$, для произвольных $x \in E,h \in \R^n$ (любое $h$ допустимо: $E$ открыто, поэтому $0$~--- внутренняя точка $I$)

  Из произвольности $h$ следует, что $\nabla^2 f(x)$  неотрицательно определено. Следует из определения \hyperref[note:771B93FA-654D-4233-9487-B4FDD356DDFE]{неотрицательно определенной симметричной матрицы}.

  А из произвольности $x \in E$, мы получаем квантор $\forall x \in E\ \nabla^2 f(x) \ge 0$ (То что требовалось в определении)

  \proofstep{$\Leftarrow$}
  Давайте возьмем произвольные $x \in E, h \in \R^{n}$ \\
  Тогда $\phi(t) := f(x + th)$  (по построению) \\
  Вторая производная все так же равна $h^T \nabla^2f(x +th) h$,
  Мы уже имеем неотрицательную определенность Гессиана,
  а следовательно вторая производная будет всегда неотрицательна
  (следует из определения неотрицательной определенности) \\
  Тогда $\phi$~--- выпуклая по \hyperref[note:9CEE854D-ED37-4C08-82EB-E41B36C43EA6]{одномерному случаю}, а $f$ по предыдущей выкладке.
\end{proof}
  % Критерий выпуклости в терминах вторых производных -- теорема 1.3

\InputIfFileExists{text/6BC7D685-C429-45A3-B563-269840CA0664-before}{}{}
% generated by roam.py handbook export -- не редактировать, править заметки
% 20260918182821-преобразования_сохраняющие_выпуклость.org (body bdaaf1ab63ff)

\begin{theorem}[{Преобразования, сохраняющие выпуклость}]\label[theorem]{note:6BC7D685-C429-45A3-B563-269840CA0664}
  \begin{enumerate}
    \item \hyperref[note:B9EF495E-CDBE-4E80-B401-8AFEAB0DCCAC]{Аффинные преобразования}: $f: E \rightarrow \R$~--- \hyperref[note:1CA14F2C-637F-41DE-9E97-72075D19E97F]{выпуклая}, $\mathcal{A}(x) = Ax + b$~--- аффинное. Тогда $g(x) = f(Ax + b)$~--- выпуклая на $\left\{ x\ |\ Ax + b \in E \right\}$
    \item $f_{\alpha}$~--- семейство \hyperref[note:1CA14F2C-637F-41DE-9E97-72075D19E97F]{выпуклых функций} на $E$, и $\sup_{\alpha} f_{\alpha}(x) < +\infty\ \forall x \in E$. Тогда $\phi(x) = \sup_{\alpha} f_{\alpha}(x)$~--- выпуклая
    \item $f(x,y): \R^n \times \R^m \rightarrow \R$~--- выпуклая по совокупности переменных $(x, y)$.
      Тогда $\phi(x) = \inf_y f(x,y)$~--- выпуклая (если $\inf$ конечен при всех $x$).
    \item $\Phi: \R^m \rightarrow \R$~--- выпуклая и неубывающая по каждому аргументу, $f_1, \dots, f_m$~--- выпуклые.
      Тогда $\Phi(f_1(x), \dots, f_m(x))$~--- выпуклая
  \end{enumerate}
\end{theorem}

\begin{proof}
  Везде $\lambda \in [0;1]$, $z = \lambda x + (1-\lambda) y$, проверяется определение \hyperref[note:1CA14F2C-637F-41DE-9E97-72075D19E97F]{выпуклой функции}.

  \proofpart{1. Аффинное преобразование}
  Пусть $Ax + b, Ay + b \in E$. \hyperref[note:B9EF495E-CDBE-4E80-B401-8AFEAB0DCCAC]{Аффинное отображение} сохраняет выпуклые комбинации: $Az + b = \lambda (Ax + b) + (1-\lambda)(Ay + b)$, так что $Az + b \in E$ (область определения $g$ выпукла) и
  \[
  g(z) = f(\lambda (Ax + b) + (1-\lambda)(Ay + b)) \le \lambda f(Ax + b) + (1-\lambda) f(Ay + b) = \lambda g(x) + (1-\lambda) g(y)
  \]

  \proofpart{2. Супремум семейства}
  Для каждого $\alpha$: $f_{\alpha}(z) \le \lambda f_{\alpha}(x) + (1-\lambda) f_{\alpha}(y) \le \lambda \phi(x) + (1-\lambda) \phi(y)$, так как $f_{\alpha} \le \phi$ поточечно.
  Правая часть не зависит от $\alpha$, поэтому можно взять $\sup$ по $\alpha$ слева: $\phi(z) \le \lambda \phi(x) + (1-\lambda) \phi(y)$.

  \proofpart{3. Инфимум по части переменных}
  Возьмем $\varepsilon > 0$ и $y_1, y_2$ такие, что $f(x, y_1) < \phi(x) + \varepsilon$, $f(y, y_2) < \phi(y) + \varepsilon$ (существуют по определению $\inf$).
  Точка $(z,\ \lambda y_1 + (1-\lambda) y_2) = \lambda (x, y_1) + (1-\lambda)(y, y_2)$, и по выпуклости $f$ по совокупности переменных
  \[
  \phi(z) \le f(z,\ \lambda y_1 + (1-\lambda) y_2) \le \lambda f(x, y_1) + (1-\lambda) f(y, y_2) < \lambda \phi(x) + (1-\lambda) \phi(y) + \varepsilon
  \]
  Устремим $\varepsilon \rightarrow 0$.

  \proofpart{4. Композиция с монотонной выпуклой функцией}
  Обозначим $F(x) = (f_1(x), \dots, f_m(x))$. По выпуклости каждой $f_k$: $F(z) \le \lambda F(x) + (1-\lambda) F(y)$ покоординатно.
  $\Phi$ не убывает по каждому аргументу, поэтому $\Phi(F(z)) \le \Phi(\lambda F(x) + (1-\lambda) F(y))$, а по выпуклости $\Phi$ правая часть $\le \lambda \Phi(F(x)) + (1-\lambda) \Phi(F(y))$.
\end{proof}
  % Преобразования, сохраняющие выпуклость -- утверждение 1.2

\InputIfFileExists{text/01-outro}{}{}
